\chapter{Research Methodology}
\label{cha:methodology}

This chapter describes the research design, dataset, and implementation approach for the two stages addressed in this individual contribution: signal denoising~(ATP-1) and data compression~(ATP-3).

\section{Research Design}
\label{sec:research_design}

This thesis follows a quantitative, comparative experimental design. For each of the two pipeline stages, multiple candidate methods spanning classical signal processing and machine learning are implemented, applied to the same dataset under identical conditions, and evaluated using objective quantitative metrics. This approach ensures that conclusions are reproducible and provide practical guidance for engineers choosing between methods in an industrial deployment.

The research was carried out in five phases:

\begin{enumerate}[leftmargin=2em]
\item \textbf{Literature review:} A review of scientific literature on signal denoising and time-series compression to identify candidate methods and establish the evaluation framework~(Chapter~\ref{cha:background}).
\item \textbf{Dataset preparation:} Loading, inspecting, and structuring the NIST Layer~01 dataset for training and evaluating all methods.
\item \textbf{Method implementation:} Coding each classical and ML method in Python with consistent function signatures, saving all outputs to CSV files and result plots.
\item \textbf{Experimental evaluation:} Running all methods on the same data splits, recording metrics, and generating comparison figures.
\item \textbf{Analysis:} Identifying the best-performing methods, interpreting the results in the context of the research questions, and documenting all implementations.
\end{enumerate}

\section{Dataset}
\label{sec:dataset}

The NIST additive manufacturing benchmark dataset~\cite{nist_dataset} is used as the primary testbed. It provides high-frequency pyrometer image sequences from a laser powder bed fusion process~(IN625, Ytterbium fibre laser at~1070\,nm, laser power~195\,W, scan speed~800\,mm/s). The per-frame maximum pixel value is extracted to produce a one-dimensional temperature time-series representing the laser melt pool temperature. Temperature is converted using the Sakuma--Hattori calibration equation supplied with the dataset:

\begin{equation}
T\,(\text{\textdegree C}) = 2.655 \cdot x_{\text{raw}} - 800.7 - 273.15
\label{eq:sakuma}
\end{equation}

For Layer~01 this yields a signal spanning approximately~373--2099\,\textdegree C with~2,065 non-zero frames. The data is split~80\,\%/20\,\% in temporal order: the first~80\,\% form the training set for ML models, the remaining~20\,\% form the test set for evaluation.

A simulated thermocouple reference signal is generated by applying a first-order low-pass filter~($\alpha = 0.08$) to the calibrated temperature signal. This reference is used as the clean training target for the CNN, LSTM, Autoencoder, and BiLSTM denoisers.

\section{Methods: Signal Denoising (ATP-1)}
\label{sec:method_denoise}

Table~\ref{tab:denoise_methods} summarises all ten denoising methods implemented and compared in this thesis.

\begin{table}[H]
\centering
\caption{All denoising methods implemented and compared in this thesis~(ATP-1).}
\label{tab:denoise_methods}
\begin{tabular}{llll}
\toprule
\textbf{Method} & \textbf{Type} & \textbf{Parameters} & \textbf{Script} \\
\midrule
Raw Baseline       & Baseline  & ---                              & --- \\
Moving Average     & Classical & k\,=\,7                          & \texttt{classical\_methods.py} \\
Median Filter      & Classical & k\,=\,7                          & \texttt{denoise.py} \\
Savitzky--Golay    & Classical & w\,=\,11, p\,=\,3                & \texttt{classical\_methods.py} \\
Gaussian Filter    & Classical & $\sigma$\,=\,3.0                 & \texttt{classical\_methods.py} \\
Kalman Filter      & Classical & Q\,=\,0.001, R\,=\,10.0         & \texttt{classical\_methods.py} \\
CNN Denoiser       & ML        & 4 Conv1d layers, W=32, 50 ep.   & \texttt{ml\_denoise.py} \\
LSTM Denoiser      & ML        & 2 layers, h=64, W=32, 50 ep.    & \texttt{ml\_denoise.py} \\
Autoencoder Den.   & ML        & bottleneck=8, W=32, 50 ep.      & \texttt{ml\_denoise.py} \\
BiLSTM Denoiser    & ML        & 2 layers, h=64, W=32, 50 ep.    & \texttt{ml\_denoise.py} \\
\bottomrule
\end{tabular}
\end{table}

\subsection{Classical Methods}

The \textbf{moving average} filter is implemented using \texttt{numpy.convolve} with a uniform kernel of width~$k=7$. The \textbf{median filter} is implemented using \texttt{scipy.signal.medfilt} with kernel size~7 and serves as the primary classical denoising method in \texttt{denoise.py}. The \textbf{Savitzky--Golay} filter is implemented using \texttt{scipy.signal.savgol\_filter} with window~11 and polynomial degree~3~\cite{savitzky1964smoothing}. The \textbf{Gaussian filter} is implemented using \texttt{scipy.ndimage.gaussian\_filter1d} with $\sigma=3.0$. The \textbf{Kalman filter} is implemented as a recursive loop with process noise~$Q=0.001$ and measurement noise~$R=10.0$.

\subsection{Machine Learning Methods}

All four ML denoisers are implemented in \texttt{ml\_denoise.py} using PyTorch~\cite{pytorch}. The input signal is normalised to~$[0,1]$ using \texttt{MinMaxScaler} and segmented into overlapping windows of length~$W=32$ with~50\,\% overlap. The median-filtered signal serves as the clean training target for all four models.

\textbf{Training settings for all ML denoisers:} batch size~32, Adam optimiser~\cite{kingma2014adam}, learning rate~$10^{-3}$, 50 epochs, MSE loss.

\subsection{ATP-1 Evaluation}

Denoising performance is measured by:
\begin{itemize}[leftmargin=2em, nosep]
\item \textbf{Noise standard deviation~(C):} standard deviation of the residual between the signal and a 20-point moving average baseline.
\item \textbf{Noise reduction~(\%):} percentage decrease in noise standard deviation relative to the raw signal.
\item \textbf{Peak preservation:} qualitative assessment of whether genuine temperature peaks are distorted.
\end{itemize}

ATP-1 is considered satisfied if the best-performing method achieves a demonstrably lower noise standard deviation than the raw signal.

\section{Methods: Data Compression (ATP-3)}
\label{sec:method_compress}

Table~\ref{tab:compress_methods} summarises all six compression methods implemented and compared in this thesis.

\begin{table}[H]
\centering
\caption{All compression methods implemented and compared in this thesis~(ATP-3).}
\label{tab:compress_methods}
\begin{tabular}{lllcc}
\toprule
\textbf{Method} & \textbf{Type} & \textbf{Parameters} & \textbf{Ratio} & \textbf{RMSE (\textdegree C)} \\
\midrule
Raw Baseline   & Baseline  & ---                         & 1.0$\times$  & 0.00  \\
Downsampling   & Classical & factor\,=\,2                & 2.0$\times$  & 21.97 \\
SVD            & Classical & rank\,=\,10, W\,=\,32       & 2.1$\times$  & 131.6 \\
Wavelet~(Haar) & Classical & 10\,\% kept                 & 5.0$\times$  & 36.55 \\
PCA            & ML        & n\,=\,3, W\,=\,32           & 10.7$\times$ & ---   \\
Autoencoder    & ML        & bottleneck\,=\,3, W\,=\,32  & 10.7$\times$ & ---   \\
\bottomrule
\end{tabular}
\end{table}

\subsection{Classical Methods}

\textbf{Downsampling} retains every other sample~(2$\times$ factor) and reconstructs via linear interpolation. \textbf{SVD} reshapes the signal into non-overlapping windows of length~$W=32$, applies \texttt{numpy.linalg.svd}, and retains~$k=10$ singular values~\cite{golub2013matrix}. \textbf{Wavelet} compression uses PyWavelets~\cite{pywavelets} to apply the Haar discrete wavelet transform, retains the top~10\,\% of coefficients by absolute magnitude, and stores only the non-zero coefficients and their indices~\cite{learnable_wavelet}. All classical methods are implemented in \texttt{compress.py}.

\subsection{Machine Learning Methods}

Both ML compressors are implemented in \texttt{ml\_compress.py}. The calibrated signal is segmented into overlapping windows of length~$W=32$ and normalised using \texttt{MinMaxScaler}.

\textbf{PCA} uses \texttt{sklearn.decomposition.PCA(n\_components=3)}~\cite{scikit-learn, jolliffe2002principal}. Compression ratio~$= 32/3 \approx 10.7\times$.

\textbf{Autoencoder} uses a PyTorch~\cite{pytorch} encoder--decoder:
\begin{itemize}[nosep]
\item Encoder: Linear(32$\to$16) $\to$ ReLU $\to$ Linear(16$\to$8) $\to$ ReLU $\to$ Linear(8$\to$3)
\item Decoder: Linear(3$\to$8) $\to$ ReLU $\to$ Linear(8$\to$16) $\to$ ReLU $\to$ Linear(16$\to$32)
\end{itemize}
Trained for 100 epochs with Adam~\cite{kingma2014adam} and MSE loss. Compression ratio~$= 32/3 \approx 10.7\times$.

\subsection{ATP-3 Evaluation}

Compression performance is measured by:
\begin{itemize}[leftmargin=2em, nosep]
\item \textbf{Compression ratio~(CR):} defined in Equation~\eqref{eq:cr}. Target: $CR > 4\times$.
\item \textbf{Reconstruction RMSE:} RMSE between the reconstructed and original calibrated signal.
\end{itemize}

ATP-3 is considered satisfied if the best-performing method achieves $CR > 4\times$ while keeping reconstruction RMSE within acceptable tolerance.

\section{Implementation Tools}
\label{sec:tools}

All methods were implemented in Python~3.10 using the following libraries:

\begin{itemize}[leftmargin=2em, nosep]
\item \textbf{NumPy} --- numerical array operations~\cite{scipy}.
\item \textbf{SciPy} --- median filter, Savitzky--Golay filter, Gaussian filter~\cite{scipy, savitzky1964smoothing}.
\item \textbf{Scikit-learn} --- PCA and data normalisation~\cite{scikit-learn}.
\item \textbf{PyTorch} --- CNN, LSTM, Autoencoder denoiser, BiLSTM, Autoencoder compressor~\cite{pytorch}.
\item \textbf{PyWavelets} --- Haar wavelet transform~\cite{pywavelets}.
\item \textbf{Matplotlib} --- comparison plots and result figures.
\end{itemize}

\section{Sources of Error and Limitations}
\label{sec:limitations-ch3}

\textbf{Simulated thermocouple training target.} The four ML denoisers are trained to reproduce the median-filtered signal rather than a physical thermocouple reference. Results may improve if a physical thermocouple is available as a clean signal target.

\textbf{Limited ML training data.} Approximately~2,065 usable NIST Layer~01 frames constrain the ML architectures to lightweight designs. Results may improve with access to multi-layer training data.

\textbf{Compression baseline dependency.} The compression stage receives the calibrated signal as input. Compression RMSE therefore reflects both the compression error and any residual calibration error from the preceding stage~(addressed in the parallel thesis by Avula Ajay Kumar).

\textbf{Classical methods not applied to PODFAM ML training.} The ML denoising models trained on NIST data cannot be directly applied to PODFAM industrial data because PODFAM has a different spatial scan format and no thermocouple reference channel for retraining. Only classical denoising and compression methods are applicable to PODFAM data without retraining.

\let\cleardoublepage\clearpage
